\chapter{Algebry Clifforda}
\label{ch:lin30-clifford}
\index{algebry Clifforda}

W algebrze zewnętrznej z Rozdziału~\ref{ch:lin11-algebry}
iloczyn wektora przez siebie jest równy zeru.
Iloczyn Clifforda uwzględnia również formę kwadratową.
Zbudujemy go jako iloraz algebry tensorowej, uzasadnimy jego
własność uniwersalną oraz dowiedziemy wzoru na wymiar
bez zakładania niezależności normalnych monomów.

\section{Algebra Clifforda: iloczyn i dowód wymiaru}
\label{sec:lin30-clifford}

Od tego miejsca $V$ jest rzeczywistą przestrzenią
euklidesową wymiaru $n$ z iloczynem skalarnym
$\langle\cdot,\cdot\rangle$ i normą
$\|v\|^2=\langle v,v\rangle$. Algebra zewnętrzna
pamięta zorientowane pola: $v\wedge v=0$.
Algebra Clifforda pamięta ponadto \emph{długość}:
kwadrat wektora jest skalarem $-\|v\|^2$.
Przyjmujemy konsekwentnie ten znak minus;
w literaturze spotyka się również konwencję przeciwną.

\begin{definicja}[Algebra Clifforda]
\index{algebra clifforda@Algebra Clifforda}
\label{def:lin30-clifford}
Niech $I_q$ będzie ideałem obustronnym algebry
tensorowej $T(V)$ z Definicji~\ref{def:lin11-tensorowa},
generowanym przez
$v\otimes v+\|v\|^2\,1$ dla wszystkich $v\in V$.
\emph{Algebra Clifforda} to
\[
 \operatorname{Cl}(V):=T(V)/I_q.
\]
Obraz $v\in V$ oznaczamy nadal $v$,
a mnożenie w ilorazie zapisujemy bez symbolu $\otimes$.
Z definicji $v^2=-\|v\|^2\,1$.
\end{definicja}

Po usunięciu z relacji składnika $\|v\|^2\,1$
otrzymujemy algebrę zewnętrzną
(Definicja~\ref{def:lin11-zewnetrzna}).
Tak więc obie algebry startują z tej samej algebry
tensorowej; zmieniamy tylko relację, przez którą
dzielimy. W algebrze Clifforda nie można zachować
zwykłego stopnia tensorowego: relacja utożsamia
wyraz stopnia $2$ ze skalarem stopnia $0$.
Pozostaje rozróżnienie stopnia \emph{parzystego}
i \emph{nieparzystego}, bo oba wyrazy relacji
mają tę samą parzystość.
Oznaczamy odpowiadające podprzestrzenie przez
$\operatorname{Cl}^0(V)$ i $\operatorname{Cl}^1(V)$.
Jest to suma prosta: ideał $I_q$ rozkłada się na części
parzystą i nieparzystą, ponieważ jego generatory są parzyste,
a każdy mnożnik po lewej i prawej można rozdzielić
według parzystości. Iloraz zachowuje więc ten rozkład.
Iloczyn słów o parzystościach $i,j\in\{0,1\}$
ma parzystość $i+j$ modulo $2$; jest to
$\Z/2\Z$-stopniowanie algebry Clifforda.

\begin{twierdzenie}[Relacja Clifforda i własność uniwersalna]
\label{thm:lin30-cliff-relacje}
Dla $v,w\in V$ w $\operatorname{Cl}(V)$ zachodzi
\begin{equation}\label{eq:lin30-cliff-relacja}
 vw+wv=-2\langle v,w\rangle\,1.
\end{equation}
Jeżeli $\mathcal A$ jest łączną algebrą rzeczywistą
z jedynką, a $f:V\to\mathcal A$ jest liniowa i spełnia
$f(v)^2=-\|v\|^2 1_{\mathcal A}$ dla każdego $v$,
to istnieje dokładnie jeden homomorfizm algebr
$\widehat f:\operatorname{Cl}(V)\to\mathcal A$
przedłużający $f$.
\end{twierdzenie}
\begin{proof}
Rozwijając $(v+w)^2$ na dwa sposoby, otrzymujemy
\[
 v^2+vw+wv+w^2=-\|v+w\|^2\,1
  =-(\|v\|^2+\|w\|^2+2\langle v,w\rangle)\,1.
\]
Po odjęciu $v^2=-\|v\|^2 1$ i
$w^2=-\|w\|^2 1$ zostaje~\eqref{eq:lin30-cliff-relacja}.
Z Twierdzenia~\ref{thm:lin11-tensorowa-uniwersalna}
odwzorowanie $f$ przedłuża się jednoznacznie do homomorfizmu
$F:T(V)\to\mathcal A$. Dla każdego generatora ideału
$I_q$ mamy
$F(v\otimes v+\|v\|^2 1)
 =f(v)^2+\|v\|^2 1_{\mathcal A}=0$.
Jądro homomorfizmu jest ideałem, więc zawiera $I_q$.
Wobec tego wzór $\widehat f([t])=F(t)$
nie zależy od przedstawiciela $t\in T(V)$
i wyznacza żądany homomorfizm na ilorazie.
Jednoznaczność wynika stąd, że $V$ i jedynka
generują $\operatorname{Cl}(V)$.
\end{proof}

Przy bazie ortonormalnej $(e_1,\ldots,e_n)$
relacja~\eqref{eq:lin30-cliff-relacja} mówi
\begin{equation}\label{eq:lin30-cliff-orto}
 e_i^2=-1,\qquad e_ie_j=-e_je_i\quad(i\ne j).
\end{equation}
Stąd każdy iloczyn generatorów można uporządkować
rosnąco, redukując powtórzone indeksy.
Trudniejszą częścią jest wykazanie, że tak otrzymane
wyrazy nie mają dodatkowych relacji.

\begin{twierdzenie}[Baza algebry Clifforda]
\label{thm:lin30-cliff-baza}
Dla bazy ortonormalnej $e_1,\ldots,e_n$ elementy
\[
 1,\qquad e_{i_1}\cdots e_{i_r}
 \quad(1\leq i_1<\cdots<i_r\leq n,\ 1\leq r\leq n)
\]
stanowią bazę $\operatorname{Cl}(V)$. Zatem
$\dim_\R\operatorname{Cl}(V)=2^n$.
Dla $n\geq1$ części parzysta i nieparzysta mają
po $2^{n-1}$ wymiarów.
\end{twierdzenie}
\begin{proof}
\emph{Rozpinanie.} Obrazy słów $v_1\cdots v_r$
rozpinają iloraz $T(V)/I_q$. Po rozwinięciu każdego
$v_j$ w bazie wystarczą słowa z liter $e_i$.
Równości~\eqref{eq:lin30-cliff-orto} pozwalają zamienić
każde dwie sąsiednie różne litery (ze zmianą znaku)
i usunąć parę równych liter ($e_i^2=-1$).
Powtarzając te operacje, dostajemy kombinację
iloczynów o rosnących, różnych indeksach.

\emph{Niezależność.} Skorzystamy z algebry zewnętrznej
$\Lambda(V)$ z Rozdziału~\ref{ch:lin11-algebry}.
Na niej określmy dwie operacje:
\[
 \varepsilon(v)\omega:=v\wedge\omega,\qquad
 \iota_v(w_1\wedge\cdots\wedge w_r)
 :=\sum_{j=1}^r(-1)^{j-1}
   \langle v,w_j\rangle\,
   w_1\wedge\cdots\wedge\widehat{w_j}
       \wedge\cdots\wedge w_r .
\]
Kreska nad $w_j$ oznacza pominięcie tego czynnika;
przyjmujemy $\iota_v(1)=0$.
Wzór na $\iota_v$ jest wieloliniowy i alternujący
w $w_j$: przy zamianie dwóch sąsiednich argumentów składniki
pomijające inny argument zmieniają znak przez antysymetrię iloczynu
zewnętrznego, a dwa pozostałe zamieniają się miejscami z przeciwnymi
znakami. Zatem własność uniwersalna
z Twierdzenia~\ref{thm:lin11-potega-wedge-uniwersalna}
zapewnia, że operator jest dobrze określony.
Bezpośrednio z tego wzoru otrzymujemy
\[
 \iota_v(w\wedge\omega)
  =\langle v,w\rangle\omega-w\wedge\iota_v\omega.
\]
Pierwszy składnik powstaje przy usunięciu pierwszego
czynnika; pozostałe mają o jeden minus więcej niż
w $\iota_v\omega$. Usunięcie dwóch czynników
w dwóch przeciwnych kolejnościach daje składniki
o przeciwnych znakach: dla pozycji $p<q$ ich znaki to
$(-1)^{p-1}(-1)^{q-2}$ oraz
$(-1)^{q-1}(-1)^{p-1}$. Po ich dodaniu
otrzymujemy zero, zatem $\iota_v^2=0$.
Również $\varepsilon(v)^2=0$, bo $v\wedge v=0$.
Jeżeli $c(v):=\varepsilon(v)-\iota_v$, to
\[
 c(v)^2
 =-\bigl(\varepsilon(v)\iota_v
          +\iota_v\varepsilon(v)\bigr)
 =-\|v\|^2\id_{\Lambda(V)}.
\]
Własność uniwersalna z poprzedniego twierdzenia
daje więc reprezentację
$c:\operatorname{Cl}(V)\to\operatorname{End}_\R(\Lambda V)$.
Przyłóżmy ją do $1\in\Lambda^0V$.
Dla różnych ortonormalnych $e_{i_j}$ wszystkie
kontrakcje pojawiające się podczas kolejnego
działania znikają: nowy $e_{i_j}$ jest prostopadły
do wcześniej utworzonych czynników. Stąd
\[
 c(e_{i_1}\cdots e_{i_r})\,1
    =e_{i_1}\wedge\cdots\wedge e_{i_r},\qquad
 c(1)\,1=1.
\]
Prawa strona przebiega bazę $\Lambda(V)$ z
Twierdzenia~\ref{thm:lin11-zew-baza}. Jeśli kombinacja
uporządkowanych słów znikałaby w algebrze Clifforda,
jej działanie na $1$ dawałoby zerową kombinację
\emph{różnych} wektorów tej bazy. Wszystkie
współczynniki muszą więc być zerowe.
Słów o $r$ indeksach jest $\binom nr$;
sumując po $r$ otrzymujemy $2^n$.
Wreszcie bijekcja $I\mapsto I\triangle\{1\}$
zamienia podzbiory parzystej i nieparzystej
liczności, co dowodzi równości ich liczby.
\end{proof}

Na przykład w $\operatorname{Cl}(\R^2)$ redukcja słowa
$e_2e_1e_2=-e_1e_2^2=e_1$ łączy zamianę sąsiednich
generatorów z relacją $e_2^2=-1$. Dzięki dowodowi
niezależności wiemy, że taka redukcja nie ukrywa
żadnych dodatkowych relacji między słowami uporządkowanymi.

\begin{uwaga}[Podobne przestrzenie, różne mnożenia]
\label{rem:lin30-cliff-zewnetrzna}
Istnieje kanoniczny izomorfizm \emph{przestrzeni liniowych}
$\Lambda V\cong\operatorname{Cl}(V)$: wysyła
$v_1\wedge\cdots\wedge v_r$ na antysymetryzację
$\frac1{r!}\sum_{\sigma\in S_r}\operatorname{sgn}(\sigma)
 v_{\sigma(1)}\cdots v_{\sigma(r)}$.
Jest dobrze określony, bo wyrażenie jest alternujące;
na bazowych wektorach ortonormalnych daje
$e_{i_1}\cdots e_{i_r}$, więc poprzednie twierdzenie
zapewnia odwracalność. Nie jest to jednak izomorfizm
algebr: w $\Lambda V$ mamy $v\wedge v=0$,
a w $\operatorname{Cl}(V)$ mamy
$v^2=-\|v\|^2$. To wyjaśnia, dlaczego obie
przestrzenie mają wymiar $2^n$, choć kodują
inne mnożenia.
\end{uwaga}

\begin{przyklad}[Wymiary $0$, $1$ i $2$]
Dla $V=0$ otrzymujemy $\operatorname{Cl}(0)=\R$.
W wymiarze $1$, przy $e_1^2=-1$, każdy element
ma postać $a+be_1$ i mnoży się jak liczba
$a+bi\in\C$. Stąd $\operatorname{Cl}(\R)\cong\C$
\emph{jako algebra rzeczywista}.
W wymiarze $2$ baza $(1,e_1,e_2,e_1e_2)$
spełnia dokładnie relacje kwaternionów:
$e_1\leftrightarrow i$, $e_2\leftrightarrow j$,
$e_1e_2\leftrightarrow k$, zatem
$\operatorname{Cl}(\R^2)\cong\mathbb H$.
Sprawdźmy brakujący kwadrat i znaki:
$(e_1e_2)^2=-e_1^2e_2^2=-1$,
$e_2e_1=-e_1e_2$,
$e_2(e_1e_2)=-e_1e_2^2=e_1$ oraz
$(e_1e_2)e_1=-e_1^2e_2=e_2$.
Relacje $i^2=j^2=-1$ i $ij=-ji$ dają homomorfizm
wysyłający $e_1,e_2$ na $i,j$.
Obrazy czterech wektorów bazy to $1,i,j,k$, czyli baza
$\mathbb H$, więc ten homomorfizm jest izomorfizmem.
Na przykład $(e_1+2e_2)^2=-5$, podczas gdy
$(e_1+2e_2)\wedge(e_1+2e_2)=0$
w algebrze zewnętrznej. Porównujemy więc kwadraty tego samego wektora.
\end{przyklad}

\begin{stwierdzenie}[Iloczyn wektorów: skalar i dwuwektor]
\label{prop:lin30-cliff-wedge}
Przez powyższy izomorfizm przestrzeni liniowych
iloczyn Clifforda dwóch wektorów rozkłada się na
\begin{equation}\label{eq:lin30-cliff-rozdzial}
 vw=-\langle v,w\rangle
       +\tfrac12(vw-wv)
 \quad\longleftrightarrow\quad
 -\langle v,w\rangle+v\wedge w.
\end{equation}
\end{stwierdzenie}
\begin{proof}
Suma $vw+wv$ jest równa
$-2\langle v,w\rangle$ z~\eqref{eq:lin30-cliff-relacja}.
Dodając połowę tej sumy do połowy różnicy,
odzyskujemy $vw$. Antysymetryzacja
$v\wedge w\mapsto\tfrac12(vw-wv)$
jest przypadkiem stopnia $2$ opisanym w
poprzedniej uwadze.
\end{proof}

Przykładowo, dla $v=e_1+2e_2$ oraz $w=2e_1+e_2$
mamy $\langle v,w\rangle=4$ i
$v\wedge w=(1\cdot1-2\cdot2)e_1\wedge e_2
=-3e_1\wedge e_2$. W algebrze Clifforda dostajemy
więc $vw=-4-3e_1e_2$; część skalarna mierzy
iloczyn skalarny, a część stopnia $2$ mierzy
zorientowane pole równoległoboku.

\begin{figure}[!htbp]
\centering
\begin{tikzpicture}[>=Stealth,line cap=round,scale=1.0]
 \coordinate (O) at (0,0);
 \coordinate (U) at (3.0,0);
 \coordinate (W) at ({3*cos(47)},{3*sin(47)});
 \coordinate (UW) at ({3+3*cos(47)},{3*sin(47)});
 \fill[manifoldgreen!16] (O)--(U)--(UW)--(W)--cycle;
 \draw[manifoldgreen!60!black,thick] (U)--(UW)--(W);
 \draw[->,manifoldblue,very thick] (O)--(U)
  node[pos=.86,below,font=\small] {$u=e_1$};
 \draw[->,manifoldred,very thick] (O)--(W)
  node[pos=.78,above left,font=\small] {$v$};
 \draw[densely dashed,manifoldred!55] (W)--({3*cos(47)},0);
 \draw[->,manifoldred!70!black] (0,-.3)--({3*cos(47)},-.3)
  node[midway,below,font=\small] {$\cos\alpha$};
 \draw[manifoldgold!80!black,thick] (.58,0)
  arc[start angle=0,end angle=47,radius=.58];
 \node[manifoldgold!80!black,font=\small] at (.83,.32) {$\alpha$};
 \node[manifoldgreen!70!black,font=\small] at (3.3,1.3)
  {$u\wedge v=\sin\alpha\,e_1\wedge e_2$};
\end{tikzpicture}
\caption{Dla jednostkowych $u=e_1$ i
$v=\cos\alpha\,e_1+\sin\alpha\,e_2$ iloczyn Clifforda
$uv=-\cos\alpha+\sin\alpha\,e_1e_2$ łączy
ujemny iloczyn skalarny ze zorientowanym polem.}
\label{fig:lin30-cliff-iloczyn}
\end{figure}

